\section{Discussion}
\label{sec:discussion}

\subsection{Key Findings and Interpretation}
\label{sec:disc_findings}

Our analysis reveals several findings with significant implications for micromobility safety policy, organized around the convergence of evidence from nine complementary models.

The most consequential finding is the overwhelming dominance of the operating mode. TUB mode increases speeding odds by 121$\times$, and the within-subject paired analysis demonstrates that the speed governor reduces speeding from 7.4\% to 0.38\% with large effect sizes ($d = 1.0$--$2.2$). The modest change in speed variability ($d = 0.22$) suggests that governors compress the speed distribution without substantially altering riding dynamics. The DiD analysis of the December 2023 TUB ban strengthens this to causal evidence: each 10~pp TUB reduction produces a 5.0~pp speeding reduction (95\% CI: [4.1, 7.4]). While the parallel trends test formally fails, the pre-trend is \textit{positive} (conservative direction), and the treatment effect is 3--4$\times$ larger than pre-period variation. This convergence across correlational, within-subject, predictive, and quasi-experimental designs provides a robust basis for policy.

A natural concern is whether former TUB riders compensate by riding faster in STD/ECO mode after losing TUB access. We address this through a user-level mode-switcher analysis (Figure~\ref{fig:mode_switcher}, Appendix). Among 37,024 ``switchers'' (TUB users pre-ban who rode STD/ECO in December) and 79,668 ``never-TUB'' riders, the within-user DiD speed change is $+0.20$~km/h ($d = 0.069$, negligible)---switchers do not accelerate their STD/ECO riding after losing TUB access. Critically, switchers' speeding rate \textit{decreased} more than never-TUB riders ($-1.71$~pp vs. $-0.27$~pp), ruling out risk compensation. The observed aggregate STD/ECO speed increase is entirely a composition effect: inherently faster riders (pre-ban gap: 0.65~km/h, $d = 0.24$) joining the STD pool raises the pool average without individual behavioral change. A placebo test comparing August--September versus October 2023 (no treatment) yields a DiD of $-0.08$~km/h ($|d| = 0.03$, negligible; Figure~\ref{fig:mode_switcher_placebo}, Appendix), confirming that the December finding is treatment-driven rather than a seasonal artifact. The convergence of null compensation at both aggregate (behavioral substitution DiD) and individual (mode-switcher DiD) levels provides strong evidence that speed governors reduce speeding without inducing compensatory behavior.

The spatial analysis confirms that speeding concentrates in specific locations (Moran's $I$: 0.16--0.71, all $p < 0.001$). The hurdle model reveals a nuanced cycling infrastructure effect: it reduces speeding probability (OR $= 0.71$) but increases intensity when speeding occurs ($\beta = 1.14$)---consistent with dedicated paths providing uninterrupted environments that both encourage moderate speeds and enable sustained speeding. This has design implications: cycling infrastructure should incorporate speed-calming features.

The four rider typologies (40\% Safe, 26\% Moderate-Risk, 26\% Habitual Speeders, 8\% Stop-and-Go) underscore heterogeneity: uniform speed reductions may be unnecessarily restrictive for the majority while insufficient for the high-risk minority. The experience analysis reveals that speeding escalation is primarily mediated by mode adoption (TUB usage: 11.5\% at trip~1 $\rightarrow$ 47.8\% at trip~50) rather than within-mode behavior changes, motivating graduated mode access policies. The curvature risk paradox---lower speeding prevalence but 5$\times$ higher risk when speeding occurs---implies dual enforcement targeting both straight (prevalent) and curvy (dangerous) segments. The SHAP analysis validates all parametric findings model-agnostically.

\subsection{Policy Implications}
\label{sec:disc_policy}

Our findings support several specific and actionable policy recommendations.

\textbf{Mandatory speed governors.} Both individual-level (20-fold speeding reduction) and population-level (28~pp reduction for high-TUB cities) evidence support mandatory speed limiting. Our data suggest the \textit{enforcement mechanism} matters more than the threshold: a 25~km/h limit with governors produces far less speeding than without.

\textbf{Geofencing.} Spatial clustering supports zone-based speed management using our hotspot maps for targeted geofence design rather than uniform city-wide restrictions.

\textbf{Infrastructure design.} Cycling infrastructure reduces speeding probability but should incorporate speed-calming features to address the conditional intensity effect on long, straight segments.

\textbf{Targeted enforcement.} The 26\% Habitual Speeder class accounts for most speeding, supporting graduated penalties over uniform restrictions.

\textbf{Graduated mode access.} Progressive TUB adoption (11.5\% $\rightarrow$ 47.8\% over 50 trips) motivates restricting new riders to governed modes, analogous to graduated driver licensing.

\subsection{Comparison with Prior Literature}
\label{sec:disc_comparison}

Our 29.6\% speeding rate is broadly consistent with \citet{cicchino2023speed}, who found 41\% of Austin riders exceeding 16~km/h, though our continuous profiles provide more precise estimates. The non-monotonic age gradient (lowest risk at 30--34) has not been previously documented for e-scooters and may reflect life-stage effects. The operating mode's explanatory power has no direct precedent---prior studies lacked within-operator mode variation. Our within-subject design substantially strengthens the between-city evidence from \citet{cicchino2023speed}. The four-class typology extends prior survey-based segmentation \citep{degele2018identifying} with objectively observed behavioral indicators.

\subsection{Limitations}
\label{sec:disc_limitations}

Several limitations should be considered when interpreting our results.

\textbf{Single operator and temporal coverage.} Data come from Swing only; other operators may differ. The primary trajectory dataset covers May 2023; seasonal variation is only partially captured by the 24-month extension (speed data available from February 2023 only).

\textbf{Measurement limitations.} The $\sim$10-second GPS interval smooths short-duration speed events. Onboard sensor speeds show strong internal consistency ($r = 1.0$ for max speed) but cannot be independently verified. We lack crash outcome data, analyzing speed as a safety proxy based on established speed-severity relationships.

\textbf{Endogeneity.} Mode choice may be endogenous (riders intending to speed select TUB), complicating causal interpretation from the within-subject analysis alone. However, the DiD analysis of the exogenous TUB ban largely addresses this concern. Korea-specific regulations and urban context may limit direct international generalizability, though the methodology is fully transferable.

\textbf{Sample selection in mode-switcher analysis.} Of approximately 240,000 users classified as switchers or never-TUB, only 117,000 (49\%) had both pre-ban and post-ban STD/ECO trips and entered the analytical sample. The approximately 50\% attrition may reflect seasonal churn, users ceasing e-scooter use entirely after the ban, or infrequent riding. If users who stopped riding were systematically different from continuers (e.g., TUB-dependent riders who had no interest in governed modes), this could introduce selection bias. However, the placebo test's null result ($d = 0.03$) suggests that the surviving sample does not exhibit differential pre-treatment trends.
